\chapter{Архитектура}

\section{Общ архитектурен модел}

Проектът използва слоеста архитектура с MVVM на ниво представяне. Това означава, че ролите на \textit{View}, \textit{ViewModel} и \textit{Model} са налице, но не са материализирани като три механични папки в един и същ проект. Вместо това те са разпределени по слоеве според отговорностите си. View частта се реализира от Blazor компонентите и Avalonia XAML екраните. ViewModel частта е концентрирана основно в приложния слой чрез споделени ViewModel-и, а при настолния клиент има допълнителен координиращ ViewModel за специфично поведение на интерфейса. Model частта е по-широко понятие и включва домейн обекти, DTO-та, заявки, service contract-и и инфраструктурни услуги.

Този подход е по-подходящ за реален .NET проект, отколкото просто физическо групиране по три папки, защото позволява домейн логиката да остане изолирана от грижите на потребителския интерфейс, а споделената приложна логика да бъде преизползвана от повече от една View технология. Диаграмата на \cref{fig:mvvm-mapping} показва точно това съответствие между MVVM ролите и конкретните части на проекта, докато \cref{fig:architecture-overview} представя по-общия слоест изглед на системата.

\optionalfigure{figures/generated/mvvm-mapping.png}{MVVM картографиране в проекта: кои части са View, ViewModel и Model}{fig:mvvm-mapping}

\optionalfigure{figures/generated/architecture-overview.png}{Обща архитектурна диаграма на проекта}{fig:architecture-overview}

Към собствените диаграми в документацията може да се добави и външен архитектурен поглед върху репото, генериран чрез \href{https://github.com/CodeBoarding/CodeBoarding}{CodeBoarding}. Тази визуализация не замества ръчния архитектурен анализ, а служи като бърза карта за ориентация в структурата на проекта и в зависимостите между основните му слоеве.

\optionalfigure[0.46\textwidth]{figures/screenshots/Codeboarding_output.png}{Външно архитектурно обобщение на \projectname, получено чрез CodeBoarding}{fig:codeboarding-output}

\section{Слоеве на системата}

\subsection{\texttt{CVAnalysisHub.Core}}

\texttt{CVAnalysisHub.Core} е домейн слоят на проекта. В него се намират основните бизнес обекти \texttt{Video}, \texttt{AnalysisRun} и \texttt{DetectionResult}. Тези типове описват бизнес данните и връзките между тях, без да зависят нито от конкретна технология за потребителски интерфейс, нито от конкретен доставчик на база данни. Така те изпълняват ролята на стабилно ядро, върху което се изграждат останалите слоеве.

\subsection{\texttt{CVAnalysisHub.Application}}

\texttt{CVAnalysisHub.Application} съдържа приложни абстракции и споделена логика на представяне. В този слой се намират DTO-тата, моделите на заявки, service contract-ите, споделените ViewModel-и, както и преизползваемите абстракции за филтри и за списъци от обекти. Именно тук е реализирана същинската идея за преизползване на ViewModel-и от различни View технологии. Вместо всяка технология да получи собствен набор от класове със сходна функционалност, и уеб, и Avalonia клиентът използват един и същ приложен слой.

\subsection{\texttt{CVAnalysisHub.Infrastructure}}

Infrastructure слоят съдържа конкретните имплементации на услугите за съхранение, медийна обработка и анализ. Тук се намират EF Core контекстът, service класовете, конфигурацията на доставчика на база данни, файловото хранилище, логиката за обработка на видео, модулите за инференция и фоновият обработчик за чакащите анализи. Това е слоят, който осигурява взаимодействието с външните технологии, без да пренася техните зависимости към домейн и приложния слой.

\subsection{\texttt{CVAnalysisHub.Web} и \texttt{CVAnalysisHub.Avalonia}}

\texttt{CVAnalysisHub.Web} представлява уеб слоя на представяне и съдържа Blazor интерфейса, преизползваеми Razor компоненти и composition root-а за уеб изпълнение. \texttt{CVAnalysisHub.Avalonia} е настолният слой на представяне и съдържа Avalonia XAML View-та, както и минимален специфичен за технологията слой за координация на настолния интерфейс. Същественото тук е, че тези два проекта изпълняват различни View роли, но не дублират основната логика на представяне, а я използват от общия приложен слой.

\subsection{Външно архитектурно резюме}

CodeBoarding дава на разработчици и coding agents визуална карта на дадена кодова база. Инструментът комбинира статичен анализ с LLM reasoning, за да генерира архитектурни диаграми, документация на ниво компоненти и навигируеми изходи, които могат да се използват в IDE, CI процеси и документация \cite{codeboarding}. В \cref{fig:codeboarding-output} той е използван като допълнителен поглед към структурата на \projectname, а не като заместител на ръчния архитектурен анализ. CodeBoarding е стартъп, започнат от Иван Милев, който е възпитаник на Технически университет -- София, продължава обучението си с магистратура в ETH Zurich и към момента развива проекта в Силициевата долина.

\begin{description}
\item[Application Logic \& State] CodeBoarding обобщава този слой като мястото, в което се оркестрира бизнес логиката и се управлява състоянието на интерфейса чрез shared ViewModel-и и service абстракции. Това съответства пряко на ролята на \texttt{CVAnalysisHub.Application}.
\item[Background Analysis Processor] Тук инструментът разпознава жизнения цикъл на задачите за анализ и координацията между медийния двигател и persistence слоя чрез фонов обработчик. Това отговаря на hosted service слоя за обработка на чакащи анализи и на свързаните processing услуги.
\item[Web Presentation Layer] Уеб клиентът е разпознат като Blazor слой на представяне, който управлява потребителските взаимодействия и браузърното изпълнение. Това е \texttt{CVAnalysisHub.Web}.
\item[Desktop Presentation Layer] Настолният клиент е разпознат като Avalonia слой на представяне, който споделя приложната логика с уеб клиента чрез общите ViewModel-и. Това е \texttt{CVAnalysisHub.Avalonia}.
\item[Computer Vision \& Media Engine] CodeBoarding отделя тежката медийна обработка, извличането на кадри и ML детекцията чрез YOLO и ONNX runtime. Това съвпада с \texttt{VideoProcessingService} и \texttt{YoloDotNetAnalysisInferenceEngine}.
\item[Data Persistence \& Infrastructure] Този слой е описан като отговорен за съхранението и извличането на данни чрез Entity Framework Core и за реализацията на service contract-ите към SQLite или PostgreSQL. Това е ролята на \texttt{CVAnalysisHub.Infrastructure}.
\item[Domain Model] Домейн моделът е разпознат като ядро от бизнес обекти, описващи основните понятия в системата, като \texttt{Video}, \texttt{AnalysisRun} и \texttt{DetectionResult}. Това съответства на \texttt{CVAnalysisHub.Core}.
\end{description}

\section{Модел на базата данни}

Моделът на съхранение е умишлено компактен. Той не е прекалено сложен, защото проектът няма нужда от голям брой таблици, а от ясна структура, която поддържа основния бизнес сценарий. Моделът е достатъчно богат, за да реализира реален работен процес за видео анализ, и в същото време остава достатъчно малък, за да бъде лесен за обяснение и демонстрация.

Схемата съдържа три основни таблици. Таблицата \texttt{Video} описва качените видеа и пази връзка към файловото хранилище. Таблицата \texttt{AnalysisRun} представя конкретно изпълнение на анализ върху дадено видео и пази информация за модела, времето на изпълнение, статуса, броя открити обекти и изходния видео резултат. Таблицата \texttt{DetectionResult} пази единичните детекции за конкретно изпълнение, включително номер на кадър, етикет, увереност и координати на ограничителната рамка. По този начин едно видео може да има множество изпълнения на анализ, а всяко изпълнение може да съдържа множество резултати от детекция.

\optionalfigure{figures/generated/db-schema.png}{ER диаграма на основните таблици в проекта}{fig:db-schema}

\section{Последователност на анализа}

Една от важните архитектурни характеристики на проекта е, че последователността на анализа не е реализирана като директно синхронно действие от слоя на потребителския интерфейс. Вместо това процесът е разделен на ясни етапи. Първо потребителят качва видео и създава задача за анализ. След това се създава \texttt{AnalysisRun} със статус \texttt{Queued}. Фоновият обработчик периодично проверява за чакащи задачи, избира следващата и я премества в състояние \texttt{Processing}. Източникът се подава към избрания модул за инференция, който извлича кадри, изпълнява детекция, създава резултатите и накрая генерира анотирано изходно видео. В зависимост от изхода статусът се променя на \texttt{Completed} или \texttt{Failed}.

Този модел прави системата по-близка до реален работен процес, отколкото до демонстрационен сценарий само за настолен клиент. Освен това той отделя потребителския интерфейс от тежката обработка и позволява по-чиста организация на отговорностите между слоя на представяне, приложния слой и инфраструктурния слой.

\optionalfigure{figures/generated/analysis-flow.png}{Диаграма на анализа на видео в проекта}{fig:analysis-flow}
